\section{Methods}\label{sec:methods}

\subsection{Substrate: from OAuth-bound API (v0.3) to Bedrock pilot (v0.4)}

The v0.3 four-arm matrix ran against the OAuth-bound Anthropic API through the \texttt{claude --print} CLI in a bytewise-isolated subprocess (subprocess flag isolation, scrubbed \texttt{HOME} that selectively symlinks the macOS keychain or \texttt{\textasciitilde/.config/claude/} for Linux, per-item \texttt{IntegrityProbe} on \texttt{(env\_hash, flags\_hash)} that halts on leakage) --- the design discharged the v0.2 Claude-Code-context contamination critique \citep{PratyakshaContextEng2025}. v0.4 keeps that substrate isolation byte-for-byte but switches the \emph{routing target} from the OAuth-bound API to AWS Bedrock, with the model identifiers \texttt{global.anthropic.claude-haiku-4-5-20251001-v1:0} for the cascade scorer and \texttt{global.anthropic.claude-sonnet-4-5-20250929-v1:0} for the judge. The substrate change was driven by quota-availability constraints during the Phase 7 pilot (the OAuth substrate had per-account session caps that would not have admitted a 1,277-call run); ADR-006 (\texttt{docs/adr/v0.4/ADR-006-haiku-rate-limit-error.md}) documents the routing target switch and the rate-limit handling. ADR-007 (\texttt{docs/adr/v0.4/ADR-007-sdk-removal.md}) documents the parallel removal of the (deprecated) Anthropic Python SDK code path: the cascade now exclusively uses the CLI substrate, which gives the architecture a single auth chain across Cursor, Claude Code, the standalone CLI, and the Bedrock pilot (\S\ref{sec:substrate}).

\subsection{Arms (the v0.3 four-arm matrix, retained)}

The four-arm matrix from v0.3 carries forward unchanged in its arm definitions, since the v0.2 confounds it was designed to control are still load-bearing:

\begin{enumerate}[leftmargin=1.4em]
\item \textbf{\texttt{haiku\_bare} (no PCE).} Claude Haiku invoked through the v0.4 isolated CLI substrate (\S\ref{sec:engine}); single sampler call; no cascade. Architecture-vs-nothing primary control.
\item \textbf{\texttt{haiku\_cascade} (with PCE).} Same Haiku endpoint through the v0.4 cascade with \texttt{commit\_policy="event\_gated"}, $K{=}3$, \texttt{max\_tokens}$=200$, the active-inference uplift (aspect-conditioned BMR, Hopfield warm-start, plumbed \iast{cit}-temperature, per-item \texttt{FreeEnergyBudget}), and an always-shadow revision that lets H8a be measurable on every item.
\item \textbf{\texttt{haiku\_bare\_2K\_scorer} (no PCE, $+K$ compute).} Single pass with \iast{icch\=a} fanning out $K' = 2K$ candidates and the same \iast{j\~n\=ana} BMR scorer the cascade uses (\texttt{commit\_policy="always\_draft"}). Spends roughly the same Haiku-token budget as the cascade but denies it the second pass and the \iast{vimar\'sa} brief. Extra-compute control (H6.v4).
\item \textbf{\texttt{haiku\_generic\_revise\_2pass} (with 2-pass, no \iast{vimar\'sa} brief).} Two passes with \texttt{commit\_policy="always\_revise"} but the \iast{vimar\'sa}-generated brief is replaced by a fixed generic creative-revise prompt (\texttt{GENERIC\_REVISE\_BRIEF}). Revision-protocol control (H7.v4).
\end{enumerate}

\subsection{Commit-policy multiplexer (the v0.4 addition)}

A v0.4 redesign of the cascade lifts the commit decision from a single configurable parameter to a \emph{multiplexer}: the cascade always runs both the draft and the shadow revision, scores both, and then commits the surface chosen by the active commit policy. Four policies are evaluated in parallel arms, each named \texttt{haiku\_cascade\_<policy>}:

\begin{enumerate}[leftmargin=1.4em]
\item \texttt{always\_draft} -- commit the draft, ignore the revision (degenerate baseline; equivalent to disabling the second pass).
\item \texttt{always\_revise} -- commit the revision unconditionally.
\item \texttt{event\_gated} -- commit the revision iff the \iast{vimar\'sa} event-disjunction fires (the existing v0.3 policy).
\item \texttt{learned\_gate} -- commit the revision iff a small logistic head trained on the H8a paired-revision data (ADR-002) predicts $P(\text{revision better})>0.5$ from a 6-feature signature of the draft (composite-score, $\Delta F$, aspect-cosine, Hopfield retrieval similarity, $|\text{revision}-\text{draft}|$ embedding distance, and length).
\end{enumerate}

The multiplexer construction means \emph{all four policies see the same draft and shadow revision} on each item; the only varying choice is which surface gets committed. This makes the H8c policy comparison a fair within-cascade contrast rather than a between-pipeline contrast. Because the same scorer ranks both surfaces, H8c also stress-tests the scorer-vs-judge agreement (H9): a policy that commits the surface the scorer prefers may not commit the surface the \emph{judge} prefers, and the H9 result tells us how often that disagreement matters.

\subsection{Pre-registered hypotheses (v0.4)}\label{sec:hypotheses}

The v0.4 hypothesis stack inherits H1--H4 directly from v0.3 (one per domain, primary cascade-vs-bare contrast), redefines H5 as a fixed-effects pool (ADR-005), retains H6 and H7 as the v0.3 fairness controls, and replaces the single v0.3 H8 with a three-part decomposition (H8a paired-revision quality, H8b gate calibration, H8c commit-policy comparison), plus a new H9 (judge-vs-scorer agreement). All hypotheses are pre-registered in \texttt{docs/SPEC\_v0.4.md} \emph{before} the Phase 7 Bedrock pilot is run:

\begin{itemize}[leftmargin=1.4em]
\item \textbf{H1.v4--H4.v4.} \texttt{haiku\_cascade} achieves higher composite than \texttt{haiku\_bare} on AUT, poetry-interp, poetry-gen, sci-creativity (one per domain).
\item \textbf{H5.v4 (re-registered, ADR-005).} The four primary contrasts are pooled via fixed-effects inverse-variance meta-analysis on per-domain Hedges' $g$. Reported as $g$, 95\% CI, and per-domain weights.
\item \textbf{H6.v4 (extra-compute fairness).} \texttt{haiku\_cascade} achieves higher per-domain composite than \texttt{haiku\_bare\_2K\_scorer}.
\item \textbf{H7.v4 (revision-protocol fairness).} \texttt{haiku\_cascade} achieves higher per-domain composite than \texttt{haiku\_generic\_revise\_2pass}.
\item \textbf{H8a.v4 (revision quality, paired).} Within \texttt{haiku\_cascade\_always\_revise}, the paired \texttt{score(surface\_revision)} $-$ \texttt{score(surface\_draft)} is positive over all items.
\item \textbf{H8b.v4 (gate calibration).} Both event-gated and learned-gate commit policies are evaluated as binary classifiers (``predict revision is better than draft''). We report precision, recall, F1, and accuracy on the H8a paired ground truth.
\item \textbf{H8c.v4 (commit-policy leaderboard).} The four commit policies are pairwise-compared on paired delta vs.\ \texttt{haiku\_bare}, with paired permutation tests across all six policy pairs.
\item \textbf{H9.v4 (judge-vs-scorer agreement).} The Sonnet-4 judge's per-item delta (cascade $-$ bare) and the Haiku composite scorer's per-item delta correlate at Spearman $\rho > 0$ (one-sided), with a sign-agreement diagnostic on the same paired items.
\end{itemize}

\subsection{Statistical protocol}

For each per-domain hypothesis we report: paired mean $\Delta$ (raw and length-controlled, where the per-arm linear word-count effect is regressed out before pairing), Hedges' $g$ with small-sample correction (raw and length-controlled), paired permutation $p$ (one-sided directional, exact for $n \le 18$, Monte-Carlo with 50,000 permutations otherwise), Wilcoxon signed-rank $p$ (one-sided), 95\% BCa bootstrap CI of the paired mean (10,000 resamples), Holm--Bonferroni adjusted $p$ across $\{H_1, H_2, H_3, H_4\}$, and a-priori (assumed $g{=}0.5$) and retrospective (observed $g$) statistical power. A hypothesis is \emph{supported} iff its Holm-adjusted $p<0.05$ \emph{and} its BCa CI is strictly positive. H5.v4 is reported as a fixed-effects pool with inverse-variance weights and a 95\% Wald CI on the pooled $g$. H8a is reported as a within-arm paired test with BCa CI. H8b is reported as binary-classifier metrics (precision, recall, F1, accuracy, support). H8c is reported as a leaderboard with paired BCa CIs and a pairwise paired-permutation matrix. H9 is reported as Spearman $\rho$, $p$, and a sign-agreement rate. Every JSON leaf passes through \texttt{\_clean\_json}: non-finite floats become \texttt{null} and the writer uses \texttt{allow\_nan=False}, so \texttt{benchmarks/results\_v0.4/stats.json} is RFC-compliant strict JSON.

\subsection{Reproducibility}

All seeds are recorded; bootstrap and permutation RNGs are seeded by \texttt{--seed} (default 4242). The benchmark driver (\texttt{benchmarks/driver.py}) writes a checkpoint after every call so a re-run resumes exactly from the failure point and runs the per-item \texttt{IntegrityProbe} (cached on \texttt{(env\_hash, flags\_hash)}; halts on leakage unless \texttt{--allow-leakage} is set). The v0.4 audit log lives under \texttt{audit/v0.4/} (cost ledgers per domain plus \texttt{audit/v0.4/cost\_ledger\_merged.json}, \texttt{audit/v0.4/integrity\_probes\_merged.jsonl}, \texttt{audit/v0.4/lit\_verification.jsonl}, \texttt{audit/v0.4/lit\_new\_entries.jsonl}, and \texttt{audit/v0.4/phase8\_preflight.txt}). Every numerical claim in this paper is generated from \texttt{benchmarks/results\_v0.4/stats.json} via \texttt{benchmarks/autoreport.py --version v0.4} and is therefore traceable end-to-end.

\subsection{TRIZ-resolved contradictions (v0.3 carried + v0.4 deltas)}

The five v0.3 TRIZ resolutions (clean substrate vs.\ OAuth, event-gated commit vs.\ measurability, BMR aspect-conditioning, Hopfield in cascade vs.\ purity, free-energy budget vs.\ open-ended cascade) are inherited unchanged. v0.4 introduces three additional ADRs:

\begin{enumerate}[leftmargin=1.4em]
\item \textbf{ADR-001 v0.4 (best-of-K \iast{cit}-temperature).} The \iast{cit}-temperature is now plumbed through the \iast{icch\=a} fan-out per-candidate, so the $K$-candidate posterior is sampled at a deliberate variance rather than at the single substrate-default temperature.
\item \textbf{ADR-002 v0.4 (learned gate).} A new commit policy that replaces the disjunctive event-gate with a logistic predictor trained on the H8a paired-revision ground truth. Tunes the gating threshold against the cascade's own measured revision quality rather than against a hand-coded $\Delta F$ heuristic.
\item \textbf{ADR-004 v0.4 (Sonnet judge protocol).} A separate judge call against \texttt{claude-sonnet-4-5-20250929-v1:0} on every cascade item, scoring the cascade-vs-bare pair under the same axis decomposition the Haiku scorer uses. The H9 agreement test is the principal output.
\end{enumerate}

Each ADR is end-to-end traceable: TRIZ card $\to$ ADR $\to$ operator code $\to$ unit test $\to$ pilot result.
